% ======================================================================
% SECTION 11: DISCUSSION
% Three thematic blocks comparing the seven proposed bills to the
% current legal stack, the AI/robotics evidence base, and the FDA's
% own RTCT framework.
% ======================================================================

[PROCESSING INSTRUCTION (Discussion block 1 - Comparison to Current
Legal Stack, 4 to 6 paragraphs): Read in full
patients/paper/Deep-Research-2/chunk\_03\_ranks10-13\_future.md (the
"Baseline versus future AI and robotics" synthesis section) and
synthesize how the seven proposed bills extend each of the layered
legal-stack ranks identified in
patients/paper/Deep-Research-2/chunk\_01\_intro\_ranks1-4.md and
chunk\_02\_ranks5-9.md.

Compose a paragraph for each of the following stack ranks, naming the
proposed bill that adapts or revises it:

(11.1.a) Rank 1, 21st Century Cures Act (PL 114-255, 2016): HR 4503
extends Cures Act broad-consent permissions into a positive procedural
modification authority; HR 4507 extends the Cures Act information-
blocking framework into a positive self-custody right. Cite
\cite{rank1_cures_act_2016} and
\cite{onc_cures_act_final_rule_page}.

(11.1.b) Rank 2, CLINICAL TREATMENT Act (PL 116-260, 2020): HR 4501
extends Medicaid coverage of routine costs into a positive self-
selection right that includes 24/7 booking and decentralized payment
parity. Cite \cite{rank2_clinical_treatment_act_2020} and
\cite{cms_ncd_3101_routine_costs}.

(11.1.c) Rank 3, ACA Section 2709 (PL 111-148, 2010): HR 4501 closes
the decentralized-component payment-parity gap that Section 2709
left unaddressed. Cite \cite{rank3_aca_clinical_trials_2010} and
\cite{usc_300gg8_clinical_trials}.

(11.1.d) Rank 4, FDORA Section 3209 (PL 117-328, 2022): HR 4506
codifies American leadership under the FDORA in silico authorization,
extending Section 3209 from animal-testing alternatives to physical-
AI executor selection. Cite \cite{rank4_fdora_fda_modernization_2022}.

(11.1.e) Rank 9, ONC HTI-1 (2023): HR 4504 extends DSI transparency
into mandatory error-rate publication and patient-initiated AI
second-opinion rights. Cite \cite{rank9_onc_hti1_2023} and
\cite{hti1_dsi_fact_sheet_2023}.

Reference the prior site documentation legislation as a state-level
parallel that informs each federal bill: SB 1042 (authorization act,
parallel to HR 4501), AB 2847 (patient rights, parallel to HR 4502
and HR 4503), and SB 892 (data protection, parallel to HR 4507).
Locations:
new-trial/site/01-legislation-authorization/physical\_ai\_trial\_authorization.tex,
new-trial/site/02-legislation-patient-rights/physical\_ai\_patient\_rights.tex,
and new-trial/site/03-legislation-data-transparency/physical\_ai\_data\_transparency.tex.
Cite \cite{kawchak_2026_19176370}.]

[PROCESSING INSTRUCTION (Discussion block 2 - Comparison to AI and
Robotics Evidence Base, 4 to 6 paragraphs): Read in full
patients/paper/Deep-Research-1/ (all three chunks) and
patients/paper/Deep-Research-3/part2\_ai\_robotics\_legislation.md and
synthesize how the seven proposed bills create the legal scaffolding
that the AI evidence base requires to mature from pilot to deployment.

Compose paragraphs covering, in order:

(11.2.a) LLM trial-matching evidence (TrialGPT, PRISM, TrialMatchAI):
the recall and accuracy figures from
patients/paper/Deep-Research-3/part2\_ai\_robotics\_legislation.md
indicate the technology is ready, but HR 4501 (self-selection) and
HR 4507 (data self-custody) are required before patients can use it.
Cite \cite{jin2024trialgpt}, \cite{gupta2024prism}, and
\cite{abdallah2026trialmatchai}.

(11.2.b) ePRO and remote symptom monitoring (Rocque 2025, ASyMS,
eRAPID): the prospective evidence supports HR 4504 and HR 4505
without overclaiming. Frame the bills as opportunistic extensions of
documented effect sizes, not as claims of new effects. Cite
\cite{rocque2025remote_symptom_monitoring}, \cite{kearney2009asyms},
\cite{maguire2017esmart}, and \cite{absolom2021erapid}.

(11.2.c) Computational pathology (Virchow): foundation-model
performance across pan-cancer and rare-cancer detection supports HR
4501 and HR 4507's same-day pathology release requirement. Cite
\cite{vorontsov2024virchow}.

(11.2.d) Bayesian urgent-care prediction (Gonzalez 2025) and AI
notification trials (Mazor 2025): the Mazor 2025 null result on AI
notifications alone is critical and informs HR 4504's "no automated-
only denial" guardrail. Cite \cite{gonzalez2025bayesian_urgent_care_nsclc}
and \cite{mazor2025ai_notifications_trial}.

(11.2.e) Robotics evidence (TARGET bronchoscopy, Cancer CARE Beyond
Walls): direct evidence supporting HR 4502 and HR 4505 patient-
choice and home-care frameworks. Cite \cite{target2025robotic_bronchoscopy}
and \cite{dronca2026cancer_care_beyond_walls}.

(11.2.f) The author's prior simulation evidence (the four LLM
simulations in new-trial/national-24-7-trial/paper/full-paper/, the
single-patient journey in patient-journey/, the 24-hour and 168-hour
sponsor simulations in sponsor/final\_paper/) supplies the
operational substrate that the bills regulate. Cite
\cite{kawchak_2026_19994945}, \cite{kawchak_2026_19119939}, and
\cite{kawchak_2026_19396256}.]

[PROCESSING INSTRUCTION (Discussion block 3 - Respectful Comparison
to FDA RTCT and Other Governing Bodies, 3 to 4 paragraphs): Read in
full new-trial/national-24-7-trial/FDA-April-2026/FDA\_RealTime\_Clinical\_Trials.md
and frame all seven proposed bills as opportunistic extensions of
the FDA's own April 28 2026 announcement, NOT as critiques of the
FDA. The discussion must be respectful throughout but opportunistic
under the patient's control based primarily on their own interests.

Cover, in order:

(11.3.a) The FDA's "demonstrated proof-of-concept" framing (TRAVERSE,
STREAM-SCLC, Paradigm Health) is the regulatory opening that HR
4505 builds on. Cite \cite{fda_rtct_press_2026} and
\cite{fda_dhts_drug_development_2026}.

(11.3.b) The FDA Oncology Center of Excellence Advancing Oncology
DCT pathway is the regulatory opening that HR 4501 builds on. Cite
\cite{fda_oce_advancing_oncology_dct_2024}.

(11.3.c) The FDA AI Regulatory Decision-Making Draft Guidance 2025
is the regulatory opening that HR 4504 and HR 4502 build on. Cite
\cite{fda_ai_regulatory_decision_draft_2025}.

(11.3.d) The HHS HIPAA Right-to-Access 2025 and OHRP Broad Consent
2017 frameworks are the regulatory openings that HR 4503 and HR 4507
build on. Cite \cite{hhs_right_to_access_2025} and
\cite{ohrp_broad_consent_attachment_c_2017}.

Close this block by stating that the United States must remain
Number 1 in the world for patient benefit during the medical AI and
robotics revolution, and that the seven bills are the legal
infrastructure that secures that leadership without rolling back
existing patient protections under 45 CFR 46, HIPAA, or FDA informed
consent guidance. Cite \cite{hhs_45cfr46_protection_human_subjects}
and \cite{fda_informed_consent_guidance_2023}.]
